\section{Erd\H{o}s Problem \#97}

\subsection*{1) ROUND-2 OBJECTIVE}

I pursue path (A): continue the proof attempt by attacking the first unresolved cardinality after Round-1. Round-1 proved that no $6$-point counterexample exists, so the next natural gap is the $7$-point case. The goal in this round is to determine the exact combinatorial structure forced on any hypothetical $7$-point counterexample.

\subsection*{2) Round-1 FOUNDATION USED}

I use the following Round-1 results without re-proving them.

\begin{itemize}
\item Round-1 theorem: no $6$-point set in the plane can have the property that every point has four other points equidistant from it. Hence any counterexample to Problem \#97 must have at least $7$ points.
\item Round-1 incidence reformulation: if $A=\{p_1,\dots,p_n\}\subset \mathbf{R}^2$ is such a configuration and $C_i$ is a circle centered at $p_i$ containing at least four other points of $A$, then
\[
S_i:=\{j\in\{1,\dots,n\}\setminus\{i\}: p_j\in C_i\}
\]
satisfies $|S_i|\ge 4$.
\end{itemize}

\subsection*{3) NEW INSIGHT / TOOL (ROUND-2)}

The new ingredient is a double-counting argument on point--circle incidences and circle--circle intersections in the first open case $n=7$. It forces exact equalities everywhere, and these equalities imply that the complements
\[
B_i:=\{1,\dots,7\}\setminus S_i
\]
form a Steiner triple system on $7$ points.

So Round-2 does not merely say ``$n=7$ is hard''; it proves that any $7$-point counterexample must carry the rigid incidence pattern of a $2$-$(7,3,1)$ design.

\subsection*{4) ATTACK PLAN (ROUND-2)}

The exact gap left by Round-1 is this: a counterexample, if one exists, has at least $7$ points, but nothing structural was known at $n=7$.

To advance strictly beyond Round-1, I prove the following sharp reduction:

\medskip

\textbf{Claim to prove.} If $A=\{p_1,\dots,p_7\}\subset \mathbf{R}^2$ has the property that for every $i$ there is a circle $C_i$ centered at $p_i$ containing at least four other points of $A$, then:
\begin{itemize}
\item each $C_i$ contains exactly four points of $A$;
\item each point of $A$ lies on exactly four of the seven circles;
\item every two circles $C_i,C_j$ meet in exactly two points of $A$;
\item the $7$ complements $B_i=\{1,\dots,7\}\setminus S_i$ are $3$-subsets with every pair of labels contained in exactly one $B_i$.
\end{itemize}

This overcomes the main Round-1 obstacle: instead of having many uncontrolled incidence patterns for $n\ge 7$, the case $n=7$ collapses to one highly constrained combinatorial model.

\subsection*{5) WORK (ROUND-2)}

\paragraph{7-point structure theorem.}
Let $A=\{p_1,\dots,p_7\}\subset \mathbf{R}^2$. Assume that for each $i\in\{1,\dots,7\}$ there exists a circle
\[
C_i=\{x\in \mathbf{R}^2:\|x-p_i\|=r_i\}
\]
centered at $p_i$ such that $C_i$ contains at least four points of $A\setminus\{p_i\}$.
Define
\[
S_i:=\{j\in\{1,\dots,7\}\setminus\{i\}:p_j\in C_i\}.
\]
Then the following hold.

\begin{enumerate}
\item $|S_i|=4$ for every $i$.
\item If
\[
t_j:=|\{\,i\in\{1,\dots,7\}: j\in S_i\,\}|
\]
is the number of circles through $p_j$, then $t_j=4$ for every $j$.
\item For every $i\neq k$,
\[
|S_i\cap S_k|=2.
\]
\item If $B_i:=\{1,\dots,7\}\setminus S_i$, then $|B_i|=3$, $i\in B_i$, and every pair of labels from $\{1,\dots,7\}$ belongs to exactly one block $B_i$.
\end{enumerate}

\paragraph{Proof.}
\emph{Step 1: no circle can contain all six other points.}
Suppose, for contradiction, that $|S_i|=6$ for some $i$. Then $C_i$ contains every point of $A\setminus\{p_i\}$.
Fix $j\neq i$. Since $|S_j|\ge 4$, the circle $C_j$ contains at least four points of $A\setminus\{p_j\}$.
Among those at most one point is $p_i$, and every other point lies on $C_i$.
Hence $C_i$ and $C_j$ have at least three common points of $A$.
But two distinct circles in the plane intersect in at most two points, and $C_i\neq C_j$ because their centers are different.
This contradiction shows that
\[
|S_i|\in\{4,5\}\qquad\text{for all }i.
\]

\medskip

\emph{Step 2: double count pairwise circle intersections through the points of $A$.}
Let
\[
m:=\sum_{i=1}^7 |S_i|=\sum_{j=1}^7 t_j.
\]
Because every $|S_i|\ge 4$, we have $m\ge 28$.

Now count triples $(j,\{i,k\})$ with $i<k$ and $j\in S_i\cap S_k$ in two ways:
\[
\sum_{1\le i<k\le 7} |S_i\cap S_k|
=
\sum_{j=1}^7 \binom{t_j}{2}.
\]
Since any two distinct circles intersect in at most two points,
\[
|S_i\cap S_k|\le 2
\quad\text{for all }i\neq k,
\]
hence
\[
\sum_{j=1}^7 \binom{t_j}{2}
=
\sum_{1\le i<k\le 7} |S_i\cap S_k|
\le 2\binom{7}{2}=42.
\]

On the other hand,
\[
\sum_{j=1}^7 \binom{t_j}{2}
=
\frac12\left(\sum_{j=1}^7 t_j^2-\sum_{j=1}^7 t_j\right)
\ge
\frac12\left(\frac{(\sum_{j=1}^7 t_j)^2}{7}-\sum_{j=1}^7 t_j\right)
=
\frac12\left(\frac{m^2}{7}-m\right),
\]
by Cauchy--Schwarz.
Since $m\ge 28$ and the function $\frac12(\frac{m^2}{7}-m)$ is increasing for $m\ge 4$,
\[
\sum_{j=1}^7 \binom{t_j}{2}
\ge
\frac12\left(\frac{28^2}{7}-28\right)=42.
\]
Therefore all inequalities above are equalities.

Equality in
\[
\sum_{j=1}^7 \binom{t_j}{2}\le 42
\]
forces
\[
|S_i\cap S_k|=2\qquad\text{for all }i\neq k.
\]
Equality in Cauchy--Schwarz forces all $t_j$ to be equal, and since their sum is $m=28$, we obtain
\[
t_1=t_2=\cdots=t_7=4.
\]
Finally,
\[
\sum_{i=1}^7 |S_i|=\sum_{j=1}^7 t_j=28,
\]
and because each $|S_i|\ge 4$, necessarily
\[
|S_i|=4\qquad\text{for all }i.
\]

\medskip

\emph{Step 3: pass to the complements.}
Define
\[
B_i:=\{1,\dots,7\}\setminus S_i.
\]
Since $|S_i|=4$, we get $|B_i|=3$. Also $i\notin S_i$, so $i\in B_i$.

For $i\neq k$,
\[
|B_i\cap B_k|
=
7-|S_i\cup S_k|
=
7-(|S_i|+|S_k|-|S_i\cap S_k|)
=
7-(4+4-2)=1.
\]
Thus any two distinct blocks $B_i,B_k$ meet in exactly one label.

Now count unordered pairs of labels contained in the blocks $B_i$.
Each block contributes
\[
\binom{3}{2}=3
\]
pairs, so altogether the $7$ blocks contribute
\[
7\binom{3}{2}=21
\]
pairs, counted with multiplicity.
But the whole label set $\{1,\dots,7\}$ has exactly
\[
\binom{7}{2}=21
\]
unordered pairs.
Since two distinct blocks intersect in only one label, no pair of labels can occur in two different blocks.
Hence every pair of labels occurs in exactly one block.

This is exactly the defining property of a Steiner triple system on $7$ points. \hfill\emph{End of proof.}

\medskip

\paragraph{Corollary.}
Any $7$-point counterexample to Problem \#97 would induce a $2$-$(7,3,1)$ design on the label set $\{1,\dots,7\}$ via the complements $B_i=\{1,\dots,7\}\setminus S_i$.

Equivalently: the first open case $n=7$ is not a free geometric search problem anymore; it is reduced to realizing a very specific incidence structure with circles centered at the points themselves.

\subsection*{6) ADVERSARIAL VERIFICATION}

I now check the new work against the main failure modes.

\begin{itemize}
\item \textbf{Did I accidentally use convexity?} No. The entire theorem is valid for arbitrary $7$-point sets in the plane. So it is genuinely stronger than the convex-polygon statement.
\item \textbf{Could two circles coincide, invalidating the intersection bound?} No. The centers are $p_i$ and $p_j$ with $i\neq j$, so $C_i\neq C_j$.
\item \textbf{Did Step 1 really rule out $|S_i|=6$?} Yes. If $C_i$ contains all six other points, then any other $C_j$ contains at least four points, at most one of which can be $p_i$; therefore at least three lie on both circles, impossible.
\item \textbf{Did the convexity inequality hide a gap?} No. I wrote it explicitly as a Cauchy--Schwarz bound on $\sum t_j^2$, and the equality case is standard: all $t_j$ are equal.
\item \textbf{Does the pair count for the blocks really imply ``exactly one''?} Yes. Pairwise block intersections have size $1$, so no unordered pair of labels can lie in two different blocks. Since the total number of pair-occurrences is exactly $\binom72$, every pair occurs once.
\item \textbf{Interaction with Round-1.} Round-1 only excluded $n\le 6$. The new theorem settles the entire combinatorial shape of the first remaining case $n=7$, so this is a strict advance.
\end{itemize}

\subsection*{7) FINAL (EXACTLY ONE)}

\textbf{UNRESOLVED (BUT STRICTLY ADVANCED).}

Round-2 proves a new theorem: every hypothetical $7$-point counterexample must realize a Steiner triple system on $7$ labels through the complements $B_i=\{1,\dots,7\}\setminus S_i$. This is substantially stronger than the Round-1 statement ``$n\ge 7$'' and isolates a sharply formulated obstruction for the first open case.

\subsection*{8) COMPLETION ESTIMATE (MANDATORY)}

\textbf{COMPLETION: 45\%}

\subsection*{9) REFERENCES}

No new external reference is used in the Round-2 theorem; the argument above is self-contained.

\bigskip
\hrule
\bigskip

